\section{Discussion}\label{sec:discussion}

\paragraph{The latency floor of a cascade.}
In a cascade the time from the end of the caller's speech to the agent's first audio is a sum: the end-of-turn decision, the
front's act and the synthesiser's first audio, each with a floor (in WaaV, Smart Turn is asked 200\,ms into a pause, an unsure
pause on a dangling word is held for 800\,ms, and the acknowledgement is due 300\,ms after the turn). End-to-end models avoid
much of this: on Full-Duplex-Bench v1 PersonaPlex reports 0.070\,s for its paper model (0.170\,s for the released checkpoint) and Moshi measures 0.265\,s~\citep{fdb_repo,roy2026personaplex},
and TML reports 0.40\,s (SR)~\citep{tml2026interaction}, a figure WaaV's design documents call out of reach for a
cascade~\citep{waav_design}; WaaV's 1.28\,s mean (median 0.68\,s; Section~\ref{sec:eval-v1}) should be read in that light. The
same leaderboard shows the other side: the fastest models take the floor in most of the user's pauses (synthetic pause TOR
0.985 for Moshi, 0.584 for PersonaPlex) while GPT-Realtime does so in 1.0\% of them~\citep{fdb_repo}; WaaV's hold rules aim at
the latter, and reach it: 2.9\% of synthetic and 9.7\% of real pauses. WaaV works on \emph{perceived} latency instead: an acknowledgement shortly after
the turn, a conversation that continues while work runs, and results delivered when they fit. For tool use the relevant
quantity is the time to a useful response; on Full-Duplex-Bench v3 the naive cascade took 10.12\,s and the fastest system
tested 4.25\,s~\citep{lin2026fdbv3}; WaaV spoke first after 2.48\,s, the soonest of all, but completed the task after
14.19\,s on average (median 7.44\,s), the slowest (Section~\ref{sec:eval-v3}). The design trades some turn latency for
control, composability and inspectability; it does not beat end-to-end models on every axis.

\paragraph{What an explicit decision buys, and what it costs.}
An explicit act is \emph{inspectable}, since every decision is logged with the observation the model saw;
\emph{constrained}, since a grammar admits only legal acts and the acknowledgement cannot assert an outcome;
\emph{separable} from the work, since the front never reports facts and the binding words come from templates and
the background's verbatim text; and \emph{replaceable}, since any model llama.cpp serves can be the front and any
model or agent the background, without retraining. The costs are real. A 2B model misreads some turns, and a
misreading has consequences: a continuer read as a request starts a task, a request read as a continuer is lost.
The reference documentation records misreadings found in live tests and the prompt changes made for them, such as
short continuers over the agent sometimes read as a \texttt{hold}~\citep{waav_imdoc}. The safeguards (interims cut
only on a sure reading, echo words are always a wait, a late decision asks the caller to repeat) add latency of
their own: a stop is cut only once the front has read the words, although the voice is ducked at the first sound,
and an answer to a question waits for a reading. The front also needs hardware that decides well within its
budget while the recogniser and synthesiser share the machine; the design's measurements found CPU-only serving
too slow for the budget it targeted, so a GPU or a fast accelerator is recommended~\citep{waav_design}.

\paragraph{Limitations.}
The release documents its limits~\citep{waav_imdoc}. Only the English language pack ships (a Hindi pack awaits review by
native speakers). The front model is required: by design there is no rule-based fallback, and the mode ends if the front's
server is not back within 150\,s. With some recognisers a long story's first comma pause ends the turn, and with a streaming
recogniser whose words arrive late the 900\,ms barge-in deadline cuts the agent before the front has read a stop; both are
recorded as known limits. The echo gate cannot separate a caller less than 3\,dB louder than the echo. Without a shared Redis
a restart loses tasks. Bud voice agents and processing graphs run their own tools and get no voice confirmation from the
gateway, and speech-to-speech providers are monitored and corrected but not prevented from speaking, except in the
half-cascade. Deployments must check the licences of the models they download; the LiveKit text turn detector, for example,
is distributed under the LiveKit Model License~\citep{livekit_turndetector}. Finally, Full-Duplex-Bench measures English
turn-taking on recorded and synthetic audio, not confirmations, truthful cancellation, echo from real loudspeakers or long
calls with several running tasks; Section~\ref{sec:case} gives a qualitative account of such a call.

\paragraph{Safety and privacy.}
Self-hosting keeps audio, transcripts and task data on the operator's infrastructure, although the speech vendors a session
names still receive the audio they process. Confirmations bind only on what was heard, sensitive values stay off the voice
channel, call text is quoted to the front as data, and background answers are checked for links, unsourced long numbers and
leaked instructions. These are safeguards, not guarantees (a paraphrased prompt leak passes the check), and operators remain
responsible for telling callers that they speak with an automated agent and for consent where calls are recorded.

\paragraph{Future work.}
The design documents name the next steps~\citep{waav_design}: reviewed language packs beyond English, starting with
Hindi; a front model fine-tuned on teacher-labelled decisions and gated on held-out sets; approvals by voice for agents
that run their own tools; and, only if an open duplex model meets the latency budget and hands work to a background
reliably, a native full-duplex regime on WaaV's own inference engine. We also plan to evaluate on TurnBench's
end-of-turn and interruption tracks~\citep{sesame2026turnbench} and on Full-Duplex-Bench v2, whose correction tasks
exercise the \texttt{update} act directly~\citep{lin2026fdbv2}.
